\chapter{High Blood Sugar}
\index{High Blood Sugar}

\section*{Definition and Overview}
High blood sugar, or hyperglycaemia, occurs when the level of glucose in the blood is too high. Glucose is the body's main fuel, and its level is normally controlled by insulin. High blood sugar most often results from diabetes, but it can also occur temporarily from illness, stress, or certain medications, and it can affect people without diabetes. Over time, persistently high blood sugar damages blood vessels and nerves, causing complications in the eyes, kidneys, heart, and feet. Early detection and treatment restore normal sugar levels and prevent complications.

\section*{Epidemiology}
High blood sugar affects around 1.2 million people with diagnosed diabetes, and a similar number has pre-diabetes, a state of elevated blood sugar that has not yet reached diabetic levels. Type 2 diabetes, the most common form, is increasing with rising rates of obesity and an ageing population. Many people with type 2 diabetes are unaware they have it for years. Gestational diabetes affects around 15\% of pregnancies. Raised blood sugar is one of the leading contributors to cardiovascular disease, kidney failure, and preventable blindness.

\section*{Aetiology and Pathophysiology}
\begin{itemize}
    \item \textbf{Type 1 diabetes:} the immune system destroys the insulin-producing cells of the pancreas
    \item \textbf{Type 2 diabetes:} insulin resistance and declining insulin production, driven by genetics, obesity, inactivity, and age
    \item \textbf{Pre-diabetes:} blood sugar above normal but below diabetic levels
    \item \textbf{Gestational diabetes:} hormone changes during pregnancy impair insulin action
    \item \textbf{Secondary causes:} medications such as corticosteroids, pancreatic disease, hormonal disorders, and illness
    \item \textbf{Pathophysiology:} without enough effective insulin, glucose cannot enter cells and accumulates in the blood; chronic hyperglycaemia damages blood vessels through inflammation and protein modification
    \item \textbf{Acute risks:} diabetic ketoacidosis in type 1 and hyperosmolar hyperglycaemia in type 2 are emergencies
\end{itemize}

\section*{Clinical Features}
\begin{itemize}
    \item Increased thirst and frequent urination
    \item Unexplained fatigue and weakness
    \item Blurred vision
    \item Weight loss despite normal or increased appetite (especially type 1)
    \item Recurrent infections, slow wound healing, and skin problems
    \item Numbness or tingling in the hands and feet from nerve damage
    \item Many people with type 2 diabetes have no symptoms at all
    \item Severe hyperglycaemia causes dehydration, drowsiness, and eventually coma
\end{itemize}

\section*{Differential Diagnosis}
\begin{itemize}
    \item Diabetes mellitus (types 1 and 2) is the main cause
    \item Pre-diabetes, diagnosed by the same blood tests
    \item Stress hyperglycaemia during acute illness, which may resolve
    \item Medication-induced hyperglycaemia, especially corticosteroids
    \item Pancreatic and hormonal diseases
    \item Type 1, type 2, and secondary causes are distinguished by clinical features and specific antibodies
\end{itemize}

\section*{When to Seek Medical Care}
Anyone with symptoms of high blood sugar, risk factors for diabetes, or a family history should be tested. People with diabetes should seek review for persistent high readings and follow their sick-day plan during illness.

\textbf{Contact your local emergency services for:}
\begin{itemize}
    \item Confusion, drowsiness, or loss of consciousness in someone with diabetes
    \item Rapid deep breathing, fruity breath, nausea, and vomiting (ketoacidosis)
    \item Very high blood sugar with severe dehydration
    \item Chest pain, breathlessness, or signs of stroke
    \item Severe abdominal pain with vomiting in a person with diabetes
\end{itemize}

\section*{Diagnosis and Investigations}
\begin{itemize}
    \item \textbf{Fasting blood glucose:} 7.0 mmol/L or higher indicates diabetes
    \item \textbf{HbA1c test:} reflects average blood sugar over the past 3 months; 6.5\% or higher indicates diabetes
    \item \textbf{Oral glucose tolerance test:} measures the response to a glucose drink
    \item \textbf{Random blood glucose:} used with symptoms
    \item \textbf{Self-monitoring:} home glucose meters track daily control
    \item \textbf{Screening:} for complications in people with diabetes, including eye, kidney, and foot checks
\end{itemize}

\section*{Management}
\begin{itemize}
    \item \textbf{Type 1 diabetes:} lifelong insulin, with education and support
    \item \textbf{Type 2 diabetes:} lifestyle change first, with metformin and other medicines added progressively
    \item \textbf{Healthy eating:} balanced meals, controlled carbohydrates, and reduced sugar
    \item \textbf{Physical activity:} at least 150 minutes weekly improves insulin sensitivity
    \item \textbf{Weight loss:} significant benefit in type 2 diabetes
    \item \textbf{Blood sugar monitoring:} self-monitoring and regular HbA1c checks
    \item \textbf{Cardiovascular risk management:} blood pressure, cholesterol, and smoking control
    \item \textbf{Complication screening:} annual eye, kidney, and foot checks
    \item \textbf{Education and support:} diabetes education programs improve outcomes
\end{itemize}

\section*{Complications}
\begin{itemize}
    \item Cardiovascular disease: heart attack, stroke, and peripheral artery disease
    \item Kidney disease progressing to kidney failure
    \item Eye disease: retinopathy, cataracts, and blindness
    \item Nerve damage: peripheral neuropathy, foot ulcers, and amputation
    \item Diabetic ketoacidosis and hyperosmolar hyperglycaemia
    \item Infections and poor wound healing
    \item Cognitive decline and depression
\end{itemize}

\section*{Prognosis and Outcomes}
With good control, people with diabetes can live long, healthy lives with a low risk of complications. Every 1\% reduction in HbA1c reduces the risk of microvascular complications by around a third. Early diagnosis of pre-diabetes allows lifestyle change that can prevent progression to diabetes altogether. Modern medicines, technology such as continuous glucose monitors, and structured care have transformed outcomes. The key predictors of good outcomes are early detection, sustained control, and regular complication screening.

\section*{Prevention and Self-Care}
\begin{itemize}
    \item Maintain a healthy weight and be physically active
    \item Eat a diet rich in vegetables, wholegrains, and lean protein, and limit sugar
    \item Have blood sugar checked if you have risk factors
    \item Take diabetes medicines and insulin exactly as prescribed
    \item Monitor blood sugar and record results
    \item Attend annual eye, kidney, and foot screening
    \item Follow a sick-day plan during illness
    \item Quit smoking and limit alcohol
\end{itemize}

\section*{Sources and Further Reading}
The following reputable sources provide further information for healthcare students and patients seeking reliable, current advice about high blood sugar.

\begin{itemize}
    \item Diabetes Australia. \textit{Information about diabetes and blood sugar.} https://www.diabetesaustralia.com.au/
    \item Healthdirect. \textit{Hyperglycaemia and high blood sugar.} https://www.healthdirect.gov.au/high-blood-sugar
    \item National Diabetes Services Scheme. \textit{Support and products for diabetes.} https://www.ndss.com.au/
    \item NHS. \textit{Hyperglycaemia: high blood sugar.} https://www.nhs.uk/conditions/high-blood-sugar-hyperglycaemia/
    \item World Health Organization. \textit{Diabetes: fact sheet.} https://www.who.int/
    \item Mayo Clinic. \textit{Hyperglycemia: symptoms and causes.} https://www.mayoclinic.org/
\end{itemize}
